\documentclass{article}

\usepackage{amsmath,amssymb}
\usepackage{xurl}
\usepackage[hidelinks]{hyperref}

% Shared commands for notes in docs/paper-gaps/.

\DeclareUnicodeCharacter{2080}{\ifmmode{}_{0}\else\textsubscript{0}\fi}
\DeclareUnicodeCharacter{2081}{\ifmmode{}_{1}\else\textsubscript{1}\fi}
\DeclareUnicodeCharacter{2082}{\ifmmode{}_{2}\else\textsubscript{2}\fi}
\DeclareUnicodeCharacter{2099}{\ifmmode{}_{n}\else\textsubscript{n}\fi}

\newcommand{\Lean}{\textsc{Lean}}
\ExplSyntaxOn
\NewDocumentCommand{\leanid}{m}{
  \ifmmode
    \textsf{\detokenize{#1}}
  \else
    \group_begin:
    \urlstyle{sf}
    \tl_set:Nn \l_tmpa_tl {#1}
    \tl_replace_all:Nnn \l_tmpa_tl {\_} {_}
    \exp_args:NV \path \l_tmpa_tl
    \group_end:
  \fi
}
\ExplSyntaxOff
\providecommand{\tr}{\operatorname{tr}}
\newcommand{\tnleanIssueBase}{https://github.com/LionSR/TNLean/issues/}
\newcommand{\tnleanPrBase}{https://github.com/LionSR/TNLean/pull/}
\newcommand{\tnleanDocBase}{https://sirui-lu.com/TNLean/docs/}
\newcommand{\ghissue}[1]{\href{\tnleanIssueBase#1}{GitHub issue \##1}}
\newcommand{\ghpr}[1]{\href{\tnleanPrBase#1}{PR \##1}}
\newcommand{\doclink}[1]{\href{\tnleanDocBase#1.html}{\path{#1}}}

% Dirac notation and the identity operator, as in blueprint/src/macros/common.tex.
% \providecommand keeps note-local definitions authoritative.
\usepackage{dsfont}
\providecommand{\ket}[1]{|#1\rangle}
\providecommand{\bra}[1]{\langle #1|}
\providecommand{\braket}[2]{\langle #1 | #2 \rangle}
\providecommand{\ketbra}[2]{|#1\rangle\!\langle #2|}
\providecommand{\Id}{\mathds{1}}
\providecommand{\one}{\mathds{1}}
\providecommand{\C}{\mathbb{C}}
\providecommand{\MN}[1]{M_{#1}(\C)}
\providecommand{\MD}{M_D(\C)}

% External referencing for paper sources.  The build helper writes sanitized
% auxiliary files under build/paper-aux/ and excludes duplicated source labels.
\usepackage{xr}

% Import a generated paper auxiliary file with a caller-supplied label prefix.
% The candidate paths support builds from docs/paper-gaps/, the repository root,
% and build/paper-aux/ itself.
\newcommand{\importpaperaux}[2]{%
  \IfFileExists{../../build/paper-aux/#2.aux}{%
    \externaldocument[#1]{../../build/paper-aux/#2}%
  }{%
    \IfFileExists{build/paper-aux/#2.aux}{%
      \externaldocument[#1]{build/paper-aux/#2}%
    }{%
      \IfFileExists{#2.aux}{%
        \externaldocument[#1]{#2}%
      }{}%
    }%
  }%
}

% General external reference macros
\newcommand{\paperref}[2]{\ref{#1-#2}}
\newcommand{\papereqref}[2]{(\ref{#1-#2})}

% CPSV16 source (1606.00608 / MPDO-22-12-17-2.tex) external document and shortcuts
\importpaperaux{cpsv16-}{1606.00608/MPDO-22-12-17-2-tnlean}
\newcommand{\cpsvref}[1]{\paperref{cpsv16}{#1}}
\newcommand{\cpsveqref}[1]{\papereqref{cpsv16}{#1}}

% Verdict marker: \gapnote{<kind>}{<status>}. Kind is the policy
% classification (clarification, local-correction, scope-restriction,
% unfaithful, false-source, open-gap); status is open, wip (elimination
% underway), resolved, or historical. Machine-read by the site index;
% prints a verdict line.
\newcommand{\gapnote}[2]{\par\noindent\textbf{Verdict:} #1 (#2).\par}


\title{Wolf Chapter~5: Abstract Multiplicative-Domain Hypotheses}
\author{}
\date{2026-07-18}

\begin{document}
\maketitle

\section{Source statement}

Let $E:M_D(\mathbb C)\to M_D(\mathbb C)$ be a complex-linear map satisfying
the Schwarz inequality
\begin{align}
  E(A^*A)-E(A^*)E(A)\succeq 0
  \qquad (A\in M_D(\mathbb C)).
  \tag{Wolf, Equation (5.2)}
\end{align}
Wolf defines
\begin{align}
  \mathcal A_R
    &=\{A:E(XA)=E(X)E(A)\text{ for every }X\},\\
  \mathcal A_L
    &=\{A:E(AX)=E(A)E(X)\text{ for every }X\},
\end{align}
and proves
\begin{align}
  A\in\mathcal A_R
    &\iff E(A^*A)=E(A^*)E(A),\\
  A\in\mathcal A_L
    &\iff E(AA^*)=E(A)E(A^*).
  \tag{Wolf, Equations (5.13)--(5.14)}
\end{align}
The local source for Equation~(5.2) is
\path{Notes/WolfNoteTexSource/ch05_schwarz_inequalities.tex},
lines~156--164; the local source for Equations~(5.11)--(5.14) is
lines~383--410 of the same file.  No Kraus representation, complete
positivity, or $2$-positivity occurs in this statement.

\section{Earlier specialization}

The existing one-sided characterizations use a unital Kraus family.
\footnote{See
\leanid{KadisonSchwarz.mem_rightMultiplicativeDomain_iff} and
\leanid{KadisonSchwarz.mem_leftMultiplicativeDomain_iff} in
\doclink{TNLean/Channel/Schwarz/MultiplicativeDomainFull}.}
They are mathematically correct specializations, but the
Kraus representation and unitality are absent from Wolf's abstract
multiplicative-domain theorem.  Their names also use the convention that the
subscript records the side on which the varying factor is appended; this
interchanges the labels ``left'' and ``right'' relative to Wolf's convention,
where the subscript records the side occupied by the fixed element.

The source-faithful formulation now uses only complex linearity and the
displayed Schwarz inequality.
\footnote{Formal declarations:
\leanid{IsSchwarzMap},
\leanid{SchwarzMap.rightMultiplicativeDomain},
\leanid{SchwarzMap.leftMultiplicativeDomain},
\leanid{SchwarzMap.mem_rightMultiplicativeDomain_iff}, and
\leanid{SchwarzMap.mem_leftMultiplicativeDomain_iff} in
\doclink{TNLean/Channel/Schwarz/AbstractMultiplicativeDomain}.}
The Kraus declarations remain available as specialized results and are no
longer used as the blueprint witness for the unrestricted source theorem.

\section{Remaining Chapter 5 boundary}

Wolf's preceding two-variable operator Schwarz inequality and its equality
criterion, involving the inverse of $E(B^*B)$ on its range, remain separate
work.  The one-variable multiplicative-domain theorem above does not assume
those results: its polarization proof applies the one-variable Schwarz
inequality directly to $A+tY$.

\section{Verdict}

The interchange of the one-sided domain names is a notational clarification:
the two conventions define the same pair of sets with opposite labels, and
their full intersections agree.  The source-faithful one-variable theorem is
complete.  The two-variable inverse-on-range inequality and its equality
criterion remain an open gap.

\bibliographystyle{alpha}
\bibliography{references}

\end{document}
